% Gustavo Lazzari -- CV Network Engineering (PT)
% Atualizado: 2026-05-25

\documentclass[10.5pt]{article}

\usepackage[margin=0.6in, a4paper]{geometry}
\setcounter{secnumdepth}{0}
\usepackage{titlesec}
\titlespacing{\section}{0pt}{5pt}{2pt}
\titlespacing{\subsection}{0pt}{3pt}{1pt}
\titlespacing{\subsubsection}{0pt}{2pt}{1pt}
\titleformat{\section}{\large\bfseries\uppercase}{}{}{}[\titlerule]
\titleformat*{\subsubsection}{\normalsize\itshape}
\usepackage{enumitem}
\setlist[itemize]{noitemsep,topsep=1pt,left=0pt .. \parindent}
\setlist[description]{noitemsep,topsep=1pt,itemsep=1pt}
\pagestyle{empty}
\pdfgentounicode=1
\setlength{\parskip}{1pt}
\linespread{1.15}
\usepackage[utf8]{inputenc}
\usepackage[T1]{fontenc}
\usepackage{hyperref}
\hypersetup{hidelinks}
\urlstyle{same}

\begin{document}

\begin{center}
    {\LARGE\bfseries Gustavo Lazzari} \\
    \vspace{2pt}
    \href{mailto:gustavotlazzari@gmail.com}{gustavotlazzari@gmail.com} $|$
    +55 41 98773-2772 $|$
    \href{https://www.linkedin.com/in/glazzari}{linkedin.com/in/glazzari} $|$
    \href{https://github.com/GLazzari1428}{github.com/GLazzari1428} $|$
    Curitiba, PR, Brasil
\end{center}

\section{Perfil}
Estudante de Ci\^{e}ncia da Computa\c{c}\~{a}o com foco em redes e infraestrutura. Projeta e opera
rede segmentada por VLANs com OPNsense, WireGuard VPN e Cloudflare Zero Trust em homelab de
produ\c{c}\~{a}o pr\'{o}pria. Cursando CCNA 200-301. Ingl\^{e}s fluente (CAE C1), alem\~{a}o B2 (DSD II).
Elegibilidade UE em jan/2027 (passaporte italiano).

\section{Forma\c{c}\~{a}o Acad\^{e}mica}
\subsection{Pontif\'{i}cia Universidade Cat\'{o}lica do Paran\'{a} $|$ {\normalfont\itshape Bacharelado em Ci\^{e}ncia da Computa\c{c}\~{a}o} \hfill 2024 -- 2028}
\begin{itemize}
    \item Conclus\~{a}o prevista: Mar\c{c}o 2028. Disciplinas relevantes: Redes de Computadores, Seguran\c{c}a de Redes, Sistemas Operacionais.
\end{itemize}
\subsection{Col\'{e}gio Su\'{i}\c{c}o-Brasileiro de Curitiba $|$ {\normalfont\itshape Diploma Bil\'{i}ngue IB} \hfill 2017 -- 2020}

\section{Experi\^{e}ncia Profissional}
\subsection{Tribunal de Justi\c{c}a do Paran\'{a} (TJPR) \hfill Curitiba, Brasil}
\subsubsection{Estagi\'{a}rio de TI \hfill Ago 2021 -- Dez 2022}
\begin{itemize}
    \item Configurou Nginx como servidor web e reverse proxy para aplica\c{c}\~{a}o PHP da Escola da Magistratura;
          deploy em AWS EC2 com automacao via GitHub Actions.
    \item Containerizou a aplica\c{c}\~{a}o com Docker e Docker Compose; administrou banco de dados MySQL.
    \item Desenvolveu backend e frontend em PHP para gest\~{a}o de certificados e m\'{o}dulos administrativos.
\end{itemize}

\section{Projetos}
\subsection{Homelab -- Arquitetura de Rede $|$ \normalfont\textit{Ambiente Pessoal} \hfill Em andamento}
\begin{itemize}
    \item Projetou e opera rede com OPNsense como firewall/roteador principal: segmenta\c{c}\~{a}o via VLANs
          (IoT isolado de produ\c{c}\~{a}o), regras de firewall por interface e NAT.
    \item Configura e mant\'{e}m WireGuard VPN para acesso remoto seguro e Cloudflare Zero Trust com
          autentica\c{c}\~{a}o TOTP (MFA) para endpoints p\'{u}blicos.
    \item Gerencia DNS recursivo/autoritativo e DHCP; filtragem de DNS via AdGuard Home; Nginx como
          reverse proxy interno com SSL/TLS.
    \item Cluster Proxmox VE de 2 n\'{o}s com VMs e containers LXC; ZFS com experi\^{e}ncia pr\'{a}tica
          em recupera\c{c}\~{a}o de falha de disco; observabilidade com Prometheus, Grafana e Scrutiny.
\end{itemize}
\subsection{Rede IoT Segmentada $|$ \normalfont\textit{ESP32 / MQTT / Frigate NVR} \hfill Em andamento}
\begin{itemize}
    \item Frota de dispositivos ESP32 em VLAN isolada, comunicando via MQTT broker local com Home Assistant;
          Frigate NVR com detec\c{c}\~{a}o de objetos por IA para vigil\^{a}ncia.
\end{itemize}

\section{Habilidades T\'{e}cnicas}
\begin{description}
    \item[Redes] OPNsense (firewall, NAT, regras), WireGuard, VLANs, DNS/DHCP, TCP/IP, Cloudflare ZT/Tunnels, Nginx, AdGuard Home
    \item[Infra/Servidores] Linux Admin, Proxmox VE, Docker/Compose, LXC, ZFS
    \item[Monitoramento] Prometheus, Grafana, Scrutiny (S.M.A.R.T.)
    \item[Dev/Ferramentas] Bash, Python (b\'{a}sico), GitHub Actions, Git, PHP, MySQL
    \item[IoT/Embarcados] ESP32, ESPHome, MQTT, Home Assistant, Frigate NVR
    \item[Idiomas] Portugu\^{e}s (Nativo), Ingl\^{e}s (Fluente -- CAE C1), Alem\~{a}o (B2 -- DSD II), Italiano (B\'{a}sico)
\end{description}

\section{Certifica\c{c}\~{o}es}
Cambridge Advanced English -- CAE C1 $|$ Deutsches Sprachdiplom -- DSD II B2 $|$ Em curso: CCNA 200-301

\end{document}
